\clearpage


\section{Realisierungsphase}
\label{sec:Realisierungsphase}
asdfasdf

\subsection{Projektstruktur}
\label{sub:Projektstruktur}
\footnotetext{Internet Explorer (bis Version 7 pflicht), Mozilla Firefox, Google Chrome, Opera.}
\begin{wrapfigure}[14]{r}[25pt]{0\textwidth}
    \tikzstyle{every node}=[draw=black,thick,anchor=west]
    \tikzstyle{new}=[dashed,fill=hig]
    \scalebox{0.8}{
        \begin{tikzpicture}
            [%
            grow via three points={one child at (0.5,-0.7) and
            two children at (0.5,-0.7) and (0.5,-1.4)},
            edge from parent path={(\tikzparentnode.south) |- (\tikzchildnode.west)}]
            \texttt{
                \node {com.<mandant>.contrax}
                child { node {party.domain}
                child { node [new] {repository}
                }
                }
                child [missing] {}
                child { node {shared.infrastructure.async}
                child { node [new] {atmosphere}
                child { node [new] {domain}
                child { node [new] {model}
                }
                child { node [new] {service}
                child { node [new] {impl}
                }
                }
                }
                }
                };
            }
        \end{tikzpicture}
    }
    \caption{Differenz zur vorherigen Packetstruktur}
    \label{app:DifferenzPacketstruktur}
\end{wrapfigure}
asdfasdf "`Domain Driven
Design"'\footnote{Definition siehe .} asdfasdf \Code{"`com.<mandant>.contrax"'} asdfasdf \Code{<mandant>}
asdfasdfasdf \AnhangK{app:Paketstruktur} asdfasdf \Abbildung{app:DifferenzPacketstruktur} asdfasdfasdf

\subsection{Einrichten/Konfigurieren von Atmosphere}
\label{sub:Einrichten/Konfigurieren von Atmosphere}
Dasdfasdf \Anhang{app:pom.xml} asdfasdf Maven\footnote{Maven ist ein "`Build-Management-Tool"' basierend auf
Java.} asdfasdf \Klasse{VmsWebAppInitializer} asdfasdf \Code{web.xml} asdfasdf \Code{"/notify/*"} asdfasdf
\Fachbegriff{Spring} asdfasdf (siehe )

\subsection{Kommunikation: Server zu Client}
\label{sub:Kommunikation: Server zu Client}
asdfasdf asdf \Anhang{app:main.js}
\begin{description}
    \item [AtmosphereHandler] asdfasdf \Klasse{AtmosphereHandler} \Code{"/notify/*"} asdfas
    \item [Broadcaster] asdfasdf \Methode{broadcast} asdfasdf \MethodenRef{AtmosphereResource\#getBroadcaster()} asdf
    aasdfasfd \sog \Code{Factory}\footnote{asdfasdf \Code{Factory}-asdfasdfe \zB asdfasdf} asdf
\end{description}

\subsubsection{Verbindungshaltung}
\label{subsub:Verbindungshaltung}
asdfasdf (\Code{\acs{SQL}HttpRequest}) asdfasdf \Fachbegriff{Resource} asdfasdf \VerweisVL{sub:Soll-Konzept} aufgezeigt.\\
asdfasdf

\subsection{Datenhaltung}
\label{sub:Datenhaltung}
asdfasdfasdf \AnhangKapitel{app:Quellcode}

\subsubsection{UserStatusRepository}
\label{subsub:UserStatusRepository}
\begin{description}
    \item [Repository]
    asdfasdf
\end{description}

Zum Zeitpunkt der Implementierung wird bisher nur der letzte Besuch der Offenen-Task-Übersicht eines jeden Benutzers im
\Fachbegriff{Repository} gespeichert. Die Ermittlung des Zeitstempels wird unten stehend erklärt. Da es sich lediglich um ein
\Code{DateTime}-Objekt handelt, könnte diese Information in der vorhandenen Datenbank gespeichert werden. Gemäß dem Abschnitt
Datenhaltung aus \VerweisVL{sub:Technisches Konzept} wurde jedoch das \Fachbegriff{Repository} entwickelt.

\subsubsection{TaskNotificationStatisticsModel}
\label{subsub:TaskNotificationStatisticsModel}
asdfasdf \VerweisVF{sub:Soll-Konzept} asdfasdf

\subsection{Datenübertragung}
\label{sub:Datenübertragung}
asdfasdf \ggfs asdfasdf

\subsection{UI-Entwurfsumsetzung}
\label{sub:UI-Entwurfsumsetzung}
asdfasdf \VerweisKL{sub:Soll-Konzept} sdfg

